\documentclass[10pt,twocolumn]{article}
\usepackage[margin=1in]{geometry}
\usepackage{amsmath,amssymb}
\usepackage{graphicx}
\usepackage{booktabs}
\usepackage{hyperref}
\usepackage{microtype}
\usepackage{enumitem}
\usepackage{caption}
\usepackage{subcaption}

\title{\textbf{Learning to Untangle: Reinforcement Learning\\for Graph Layout Optimization}}
\author{Shenqian Wen, Zhuocheng Yu}

\date{}

\begin{document}
\maketitle

% ─────────────────────────────────────────────────────────────────────────────
\begin{abstract}
Graph drawing aims to produce visually clear 2D layouts of graphs, with
minimizing edge crossings as a key aesthetic criterion.  This work frames
layout optimization as a Markov Decision Process and trains a Proximal Policy
Optimization (PPO) agent to iteratively reposition nodes.  We design and
evaluate five progressively refined approaches, starting from a discrete-action
GAT policy that moves nodes in eight compass directions, progressing through
continuous-action and sequential placement formulations, and culminating in a
GAT-based iterative refinement policy that cycles through nodes in BFS order
with continuous delta moves.  All methods are trained and evaluated on the
Rome graph benchmark against a \texttt{neato} force-directed baseline.
Our best method, Sequential Refinement, achieves an SPC of $-24.43\%$
relative to \texttt{neato}, demonstrating that iterative local refinement
from a strong initialization substantially outperforms both from-scratch
placement and global discrete-action approaches.
\end{abstract}

% ─────────────────────────────────────────────────────────────────────────────
\section{Introduction}
\label{sec:intro}

Graph visualization is essential in network analysis, software engineering,
and biology.  Among the many quality criteria for a graph drawing, the number
of edge crossings is one of the most studied: fewer crossings reduce visual
clutter and make the underlying structure easier to perceive.

Traditional approaches such as force-directed algorithms (Fruchterman--Reingold,
Kamada--Kawai) minimize a global energy function and are computationally
efficient, but they are susceptible to poor local optima and cannot easily
incorporate combinatorial objectives like crossing minimization directly.
The crossing count is a piecewise-constant, non-differentiable function of
node positions, which means gradient-based methods can only optimize
differentiable proxies rather than the metric itself.

Reinforcement learning (RL) offers a natural alternative.  An agent interacts
with the layout environment by moving nodes, and the reward directly reflects
the hard crossing count---no differentiable surrogate is needed.  The policy
can explore large layout spaces and learn to prioritize moves that reduce
crossings, something that is difficult to express as an energy function.

This report investigates how to best formulate the RL problem for graph layout
and how increasingly expressive policy designs affect crossing reduction.
Our main contributions are:
\begin{enumerate}[leftmargin=*,topsep=2pt,itemsep=0pt]
    \item A systematic study of five RL formulations for crossing minimization,
          ranging from discrete coarse-grained actions to continuous adaptive
          node selection.
    \item \textsc{SequentialRefinement} as the best-performing method,
          achieving SPC $-24.43\%$ vs.\ \texttt{neato} by iteratively
          refining a force-directed initialization with BFS-ordered delta moves.
    \item Evaluation on the Rome graph dataset with comparisons across all
          five methods and the \texttt{neato} force-directed baseline.
\end{enumerate}

% ─────────────────────────────────────────────────────────────────────────────
\section{Background}
\label{sec:background}

\subsection{Graph Drawing and Edge Crossings}

Given an undirected graph $G=(V,E)$, a 2D layout assigns each node a
coordinate in the plane.  Two edges \emph{cross} if their straight-line
segments intersect at an interior point.  The total crossing count
$\text{cr}(\phi)$ measures the visual complexity of layout $\phi$.
Finding a minimum-crossing layout is NP-hard in general, and the best known
exact algorithms are only practical for very small graphs.

\subsection{Proximal Policy Optimization}

PPO~\cite{schulman2017ppo} is an actor-critic RL algorithm.  It collects
trajectories under the current policy, computes Generalized Advantage
Estimates (GAE), and updates the policy by maximizing a clipped surrogate
objective that prevents large policy steps.  The value function is trained
simultaneously with a coefficient of~0.5.

\subsection{Graph Attention Networks}

Graph Attention Networks (GAT)~\cite{velickovic2018gat} compute node
embeddings by aggregating neighbor information with learned attention weights.
We use three GAT layers with four attention heads and hidden dimension~128,
incorporating edge features (current Euclidean edge length) via the
edge-feature variant in PyTorch Geometric.  A global mean-pooling layer
produces the scalar state-value estimate for the critic head.

% ─────────────────────────────────────────────────────────────────────────────
\section{Problem Formulation}
\label{sec:problem}

\paragraph{State.}
All methods represent the layout as per-node feature vectors containing
normalized $(x, y)$ coordinates in $[-1,1]^2$; the remaining features differ
by method and are described in Section~\ref{sec:methods}.
Where applicable, edge features carry the current Euclidean distance between
the two endpoint nodes.

\paragraph{Reward.}
All methods share the same high-level reward structure: a crossing term
weighted by $w_c=1.0$ plus a structure term weighted by $w_s=0.3$, with
an optional $+10$ bonus for reaching a crossing-free layout.  The details
differ across methods and are described in Section~\ref{sec:methods}.
The two structure losses used are: (1)~\textbf{SoftmaxRanking loss} ---
the KL divergence between softmax distributions of layout distances and
graph-theoretic shortest-path distances (Methods 1 and 2); and
(2)~\textbf{Kamada--Kawai stress} --- a weighted sum of squared differences
between layout and graph distances (Methods 3, 4, and 5).  Method~3 computes
stress over the \emph{placed subgraph only} (the $k$ nodes placed so far),
with a full $O(k^2)$ recomputation each step; Methods 4 and~5 compute stress
over the \emph{full graph} using a precomputed weight matrix and an $O(n)$
incremental update that only touches terms involving the moved node.

During training, crossing detection can optionally use a sigmoid-based soft
approximation that provides a denser gradient signal.  Hard integer counts
are always used for evaluation and episode termination.

\paragraph{Termination.}
An episode ends successfully when all crossings are eliminated.  It is
truncated when the step budget is exhausted, when the agent's displacement
falls below a minimum threshold (indicating the layout is stuck), or when
a \emph{patience} counter---tracking consecutive steps without crossing
improvement---exceeds a threshold.  In the refinement variants, patience
can trigger a \emph{reset-to-best} mechanism: rather than truncating, the
layout is restored to the best seen so far and node traversal order is
reshuffled, giving the agent a fresh start from a good configuration.

% ─────────────────────────────────────────────────────────────────────────────
\section{Methods}
\label{sec:methods}

We develop five methods in order of increasing sophistication.  Each is
motivated by a specific limitation of its predecessor.

\subsection{Method 1: Discrete PPO}

The first formulation uses a discrete action space: at each step, the agent
selects one of $n \times 8 \times 3$ actions, encoding a choice of node,
direction (8 compass directions), and step scale (3 sizes linearly spaced
between 0.05 and 0.35 in normalized coordinates).  This multi-scale design
allows both coarse exploratory moves and fine-grained adjustments without
committing to a fixed step size.

Node features are $[x,\, y,\, \text{degree},\, \text{in\_crossing}]$ (4-dim);
edge features are $[\text{edge\_length}]$.
The policy (\textsc{DiscreteGNNPolicy}) encodes these through three GAT
layers, produces per-node logits for all direction-scale combinations,
flattens them into a single Categorical distribution, and samples one action.
The critic head uses global mean pooling to estimate the state value.
The crossing reward is normalized by the initial crossing count; the structure
term uses SoftmaxRanking loss.  A $+10$ bonus is given upon reaching zero
crossings.

The main limitation of this approach is that the action space grows linearly
with the number of nodes, and the policy must simultaneously decide both
\emph{which} node to move and \emph{how} to move it.  At the start of
training, all nodes are equally likely to be selected, which wastes many
steps on nodes that are not involved in any crossings.

\subsection{Method 2: Continuous PPO}

The second method uses a continuous action space: at each step the policy
outputs a displacement matrix of shape $(n \times 2)$, moving \emph{all}
nodes simultaneously by their respective $(\Delta x, \Delta y)$ deltas.
Node features are $[x,\, y,\, \text{degree}]$ (3-dim); edge features are
$[\text{edge\_length}]$.  The crossing reward is normalized by the maximum
possible crossing count (rather than the initial count), and the structure
term uses SoftmaxRanking loss, same as Method~1.  A $+10$ bonus is given
upon reaching zero crossings.

Moving all nodes at once gives the policy global reach each step, but it
also makes the optimization landscape more complex: a displacement that
helps one crossing cluster may worsen another.  The lack of node-level
selectivity means the policy cannot easily focus its updates on the most
problematic regions of the layout.

\subsection{Method 3: Sequential PPO}

The third method uses a Transformer-based policy that places nodes one by one
in BFS order.  We explored two variants.  In the \emph{from-scratch} variant,
the agent outputs absolute coordinates $(x, y) \in [-1,1]^2$ for each node
with no initial layout; in the \emph{neato+offset} variant, nodes start at
their \texttt{neato} positions and the agent outputs a delta offset
$(\Delta x, \Delta y)$.  Both use the same Transformer encoder, which attends
over all previously placed nodes to inform the next placement.

The from-scratch variant proves extremely difficult to train: it achieves SPC
in the range of $+80\%$ to $+90\%$, meaning the resulting layouts have far
more crossings than the \texttt{neato} baseline.  Without any initialization,
the agent must solve a global combinatorial problem from an empty canvas, and
errors made early in the BFS sequence compound through all subsequent placements.
The neato+offset variant substantially reduces this burden and achieves SPC
$+1.11\%$---still slightly worse than \texttt{neato}, but far better than
from-scratch.  The reported result in Table~\ref{tab:results} is for the neato+offset
variant.  The observation at each step is the placed subgraph: each already-placed
node contributes features $[x,\, y,\, \text{is\_neighbor\_of\_}v_t]$ (3-dim),
and edge features are $[\text{edge\_length}]$ for edges within the placed
subgraph.  The reward at each step is the negative number of new crossings
introduced by placing node $v_t$, minus a structure penalty computed as the
change in Kamada--Kawai stress over the placed subgraph (the $k$ nodes placed
so far), fully recomputed in $O(k^2)$ each step.  There is no crossing
normalization and no terminal bonus.

\subsection{Method 4: Sequential Refinement}

The fourth method returns to a refinement paradigm but uses a stronger
initialization: all episodes start from a \texttt{neato} force-directed
layout.  The agent then applies continuous delta moves $(\Delta x, \Delta y)$
to one node per step, cycling through all nodes in BFS order from a
randomly chosen start node.

The policy (\textsc{GATRefinementPolicy}) encodes a 5-dim node feature
vector $[x,\; y,\; \textit{is\_vt},\; \textit{is\_nbr},\; \textit{in\_xing}]$
through three GAT layers, where \textit{is\_vt} flags the current target
node, \textit{is\_nbr} flags its neighbors, and \textit{in\_xing} flags
crossing nodes.  It extracts the target node's embedding and outputs a
Gaussian distribution over 2D displacements, keeping the action dimension
fixed at~2 regardless of graph size.  Edge features are $[\textit{edge\_length}]$.

The reward uses the raw crossing delta (unnormalized) and Kamada-Kawai
stress as the structure term, updated incrementally in $O(n)$ per step.
There is no $+10$ terminal bonus.  Starting from a good initialization
dramatically reduces the number of steps needed to reach a low-crossing
layout.  The remaining limitation is that BFS cycling still visits all
nodes in a fixed order, regardless of whether they are involved in crossings.

\subsection{Method 5: Node-Select Refinement}

The fifth method addresses the node-selection problem directly.  Instead of
cycling nodes by an external schedule, the policy itself decides which node
to move at every step.  The action is a matrix of shape $n \times 2$; the
environment identifies the node with the largest output magnitude as the
chosen target and applies its displacement.

The policy (\textsc{Node\-Select\-GNN\-Policy}) encodes a 5-dim node feature
vector $[x,\; y,\; \textit{xing\_cnt},\; \textit{xing\_nbr},\; \textit{in\_xing}]$,
where \textit{xing\_cnt} is the normalized per-node crossing count,
\textit{xing\_nbr} flags nodes adjacent to a crossing node, and
\textit{in\_xing} flags crossing nodes; edge features are $[\textit{edge\_length}]$.
It extends the GAT refinement architecture with two new components, as follows.  First, a \emph{node selection head}
scores each node by combining a learned linear projection of its GAT
embedding with a crossing-count heuristic: nodes involved in more crossings
receive a higher base score.  The weight of the heuristic term is itself a
learnable scalar $\hat\alpha$, initialized at~2.0, allowing the model to
calibrate how much to trust the heuristic signal relative to what it learns
from experience.

Second, a \emph{crossing-context pooling} module computes a global summary
of the embeddings of all nodes currently in crossings.  This context vector
is concatenated with the selected node's embedding before the displacement
prediction head, so the predicted move is conditioned on the global crossing
structure, not just the local neighborhood.

The log-probability of a full action decomposes into the log-probability of
the node selection (a Categorical over $n$ nodes) plus the log-probability
of the displacement given the selected node.  Both terms flow through the
PPO objective, so the policy is trained end-to-end to select good nodes and
move them well.

\begin{table*}[t]
\centering
\caption{Summary of the five methods.}
\label{tab:methods}
\small
\setlength{\tabcolsep}{4pt}
\begin{tabular}{llll}
\toprule
Method & Action type & Node selection & Policy \\
\midrule
Discrete PPO        & Discrete $n{\times}8{\times}3$ & Joint w/ move & GAT \\
Continuous PPO      & Continuous 2D & All (parallel) & GAT \\
Sequential PPO      & Continuous 2D & BFS order      & Transformer \\
\textbf{Seq.\ Refinement} & Continuous 2D & \textbf{BFS shuffle} & \textbf{GAT} \\
Node-Select         & Continuous 2D & Learned        & GAT+ctx \\
\bottomrule
\end{tabular}
\end{table*}

% ─────────────────────────────────────────────────────────────────────────────
\section{Experimental Setup}
\label{sec:experiments}

\subsection{Dataset}

We use the \textbf{Rome} graph collection~\cite{didimo2011rome}, a standard
benchmark in graph drawing.  The collection contains real-world graphs from
circuit schematics, biological networks, and social graphs, split into a
training set and a held-out test set (file lists in
\texttt{data/train\_graph.txt} and \texttt{data/test\_graph.txt}).
For each graph, we precompute the edge index, per-node degree, all-pairs
shortest-path distances, and a precomputed SoftmaxRanking distance
distribution used by the structure loss.  A new graph is sampled from the
training set at the start of each episode, so the agent must generalize
across different graph topologies.

\subsection{Baselines}

The primary baseline is the \texttt{neato} layout from Graphviz, which
implements the Kamada--Kawai stress-minimization heuristic.  We record the
crossing count of the \texttt{neato} layout before any RL refinement.  All
refinement methods (Methods 4 and~5) are initialized from this layout, so
the improvement numbers directly measure how much the RL agent contributes
on top of the force-directed result.

\subsection{Training Details}

Training budgets vary by method: simpler methods (Discrete PPO, Node-Select
Refinement) are trained for $5\times10^5$ timesteps, while Sequential PPO
and Sequential Refinement use $3\times10^6$ and $2\times10^6$ timesteps
respectively due to their harder optimization landscapes.
Table~\ref{tab:hparams} lists the PPO hyperparameters for Sequential
Refinement, the best-performing method.  Evaluation is performed on the
test split with the final checkpoint.

\begin{table}[h]
\centering
\caption{PPO hyperparameters for Sequential Refinement (best method).}
\label{tab:hparams}
\small
\begin{tabular}{lc}
\toprule
Hyperparameter & Value \\
\midrule
Learning rate & $3\times10^{-4}$ \\
Discount $\gamma$ & $0.99$ \\
GAE $\lambda$ & $0.95$ \\
Clip $\epsilon$ & $0.2$ \\
PPO update epochs & $4$ \\
Batch size & $512$ \\
Rollout steps & $512$ \\
Entropy coefficient & $0.05$ \\
Hidden dimension & $128$ \\
GAT layers / heads & $3$ / $4$ \\
Structure weight $w_s$ & $0.3$ \\
Patience & $100$ steps \\
Reset-to-best & enabled \\
\bottomrule
\end{tabular}
\end{table}

\subsection{Evaluation Metrics}

Our primary metric is the \textbf{Symmetric Percentage Change (SPC)},
defined as
\[
  \text{SPC} = \frac{\text{cr}_{\text{ours}} - \text{cr}_{\text{ref}}}
                    {\max(\text{cr}_{\text{ours}},\,\text{cr}_{\text{ref}})},
\]
where $\text{cr}_{\text{ours}}$ is the crossing count produced by our method
and $\text{cr}_{\text{ref}}$ is that of the reference layout (\texttt{neato}).
SPC lies in $[-1, 1]$; a negative value means our method produces fewer
crossings than the reference.  The symmetric denominator ensures the metric
is bounded and comparable across graphs with very different crossing counts.
We also report the \textbf{zero-crossing rate}, the fraction of test graphs
for which the agent achieves a completely crossing-free layout.

% ─────────────────────────────────────────────────────────────────────────────
\section{Results}
\label{sec:results}

\subsection{Quantitative Comparison}

Table~\ref{tab:results} summarizes test-set performance across all methods.
The \texttt{neato} baseline already produces low-crossing layouts for many
small, sparse graphs, but struggles on denser graphs where the energy
landscape has many local minima.

Discrete PPO (Method~1) and Continuous PPO (Method~2) both refine a
\texttt{neato} initialization but differ in action space.  Continuous PPO
achieves SPC $-2.38\%$ and mean best-crossing count 28.65, outperforming
Discrete PPO ($-0.87\%$, mean best-cr 28.93) because it can express precise
displacements without a discretized direction grid.  However, Continuous PPO
exhibits severe instability: its mean final-state crossing count is~1{,}200,
far above its best-seen count of 28.65, indicating that the policy continues
to degrade the layout after reaching its best configuration.

The BFS placement approach (Method~3) is the only one that performs
\emph{worse} than \texttt{neato} (SPC $+1.11\%$, mean best-cr 29.25).
Even with a \texttt{neato} initialization, BFS sequential placement still
falls slightly short: each node is placed in turn and the Transformer must
predict a good position using only the already-placed subgraph, so
early placement errors compound through all subsequent nodes and the
agent cannot revise earlier decisions.

The two refinement methods (4 and~5) substantially outperform Methods 1--3,
confirming that a \texttt{neato} initialization provides a strong advantage.  Sequential Refinement
(Method~4) achieves the best overall result: SPC $-24.43\%$ and mean
crossing count 22.30, a reduction of roughly 7~crossings per graph on
average.  Node-Select Refinement (Method~5) achieves SPC $-5.20\%$,
significantly worse despite its more sophisticated architecture.  We
analyze this gap in the ablation section below.

\begin{table*}[h]
\centering
\caption{Test-set results. SPC $<0$ means fewer crossings than \texttt{neato}.}
\label{tab:results}
\small
\begin{tabular}{lccc}
\toprule
Method & Mean best-cr $\downarrow$ & SPC $\downarrow$ & Zero-cr (\%) $\uparrow$ \\
\midrule
\texttt{neato} (init)             & ---   & 0.0\%     & ---  \\
Discrete PPO                      & 28.93 & $-0.87\%$ & 3.0  \\
Continuous PPO                    & 28.65 & $-2.38\%$ & 4.0  \\
Sequential PPO                    & 29.25 & $+1.11\%$ & 3.0  \\
\textbf{Seq.\ Refinement}         & \textbf{22.30} & $\mathbf{-24.43\%}$ & \textbf{4.0} \\
Node-Select Refine                & 27.43 & $-5.20\%$ & 4.0  \\
\bottomrule
\end{tabular}
\end{table*}

\subsection{Ablation Studies}

\paragraph{Patience and reset-to-best.}
Without the patience mechanism, the agent tends to oscillate once it reaches
a local crossing minimum---it keeps moving the same nodes back and forth
without making progress.  Adding a patience counter of~100 steps truncates
these unproductive episodes early and improves sample efficiency.  Enabling
the \texttt{reset\_to\_best} option yields a further improvement: when the
patience threshold is reached, the layout is restored to the best crossing
count seen so far and the node order is reshuffled, giving the agent a fresh
perspective from a good starting point rather than terminating the episode.

\paragraph{Soft vs.\ hard crossing reward.}
The hard crossing count is piecewise-constant and provides a sparse reward
signal: it only changes when an actual crossing appears or disappears.  The
soft sigmoid approximation provides a denser signal that changes smoothly
with node positions, which helps early in training.  In practice, the two
variants converge to similar final crossing counts, suggesting that the hard
count is a sufficient training signal once the policy has learned to reliably
move nodes in the right direction.

\paragraph{Node traversal order in Sequential Refinement.}
Within the Sequential Refinement framework we compare three traversal
strategies: BFS from a random start, ascending degree (\texttt{degree\_asc}),
and descending degree (\texttt{degree\_desc}).  Table~\ref{tab:order} shows
that BFS achieves the lowest crossing count on all metrics.  Degree-descending
order performs worst, likely because moving high-degree nodes early causes
large structural disruptions that later low-degree nodes cannot repair.
BFS visits nodes in connectivity order, keeping related nodes coherent and
enabling more targeted local adjustments.

\begin{table}[h]
\centering
\caption{Node traversal order ablation (Sequential Refinement).}
\label{tab:order}
\footnotesize
\setlength{\tabcolsep}{4pt}
\begin{tabular}{lcccc}
\toprule
Order & SPC $\downarrow$ & Mean cr & Median & Max \\
\midrule
\textbf{BFS}  & $\mathbf{-24.43\%}$ & \textbf{22.30} & \textbf{8.0} & \textbf{171} \\
degree\_asc   & $-22.44\%$          & 23.28          & 8.0          & 173          \\
degree\_desc  & $-20.55\%$          & 23.51          & 10.0         & 184          \\
\bottomrule
\end{tabular}
\end{table}

\subsection{Qualitative Analysis}

Figure~\ref{fig:layouts} shows a representative example from Sequential
Refinement on graph \texttt{grafo10005.39}.  Starting from 23 crossings
in the \texttt{neato} initial layout (left), the agent reduces crossings
to 15 at episode end (center, $\Delta=-8$) and reaches a best of 13
during the episode (right, $\Delta=-10$).  Crossing edges are highlighted
in red; the rightmost panel shows how the agent consolidates nodes into a
more planar configuration by iteratively resolving the densest crossing
clusters.

\begin{figure}[h]
\centering
\includegraphics[width=\columnwidth]{seq_refinement_example.png}
\caption{Sequential Refinement on \texttt{grafo10005.39}.
  Left: initial layout (23 crossings).
  Center: final layout (15 crossings, $\Delta=-8$).
  Right: best layout seen during the episode (13 crossings, $\Delta=-10$).
  Red edges are crossing pairs.}
\label{fig:layouts}
\end{figure}

% ─────────────────────────────────────────────────────────────────────────────
\section{Discussion}
\label{sec:discussion}

\paragraph{Why RL suits crossing minimization.}
The crossing count is a combinatorial metric that is non-differentiable with
respect to node positions.  Gradient-based optimizers must resort to
differentiable surrogates, which may not correlate perfectly with the true
crossing count.  RL sidesteps this issue entirely: the policy is optimized
against the hard metric through interaction, and any proxy (such as our soft
crossing reward) is used only to smooth early training rather than as the
final objective.

\paragraph{Action space design.}
Moving from a discrete to a continuous action space (Methods 1 to 2) removes
the arbitrary direction grid and allows more precise moves.  The large
performance gap between the BFS placement approach (Method~3)
and the refinement-based approaches (Methods 4--5) confirms that initialization
quality is the single most important factor: even a simple BFS cycling strategy
on top of a \texttt{neato} layout beats from-scratch placement by a wide margin.

\paragraph{Limitations.}
Three limitations stand out from the experimental results.

First, all refinement methods depend entirely on the \texttt{neato}
force-directed initialization.  The RL agent is a \emph{post-hoc corrector}
rather than an independent layout algorithm: if \texttt{neato} produces a
layout with many crossings, the agent has more room to improve, but it
inherits all structural decisions (relative node positions, overall scale)
from the baseline.  A two-stage pipeline---first generating a coarse layout,
then refining it---or a curriculum that trains the agent on layouts of
increasing quality, could reduce this dependency.

Second, training budgets are not uniform across methods: simpler methods
receive $5\times10^5$ timesteps while Sequential PPO and Sequential
Refinement receive $3\times10^6$ and $2\times10^6$ respectively.  This
disparity makes direct cross-method comparisons difficult to interpret---a
weaker method given more compute might close the gap, and a stronger method
given less might not fully converge.  Controlled experiments with a fixed
budget would be needed to isolate architectural differences from training
effort.

Third, the attempt to replace fixed BFS cycling with a learned node-selection
policy (Method~5, Node-Select Refinement) largely failed: it achieves only
SPC $-5.20\%$ versus $-24.43\%$ for plain BFS cycling.  The root cause is
a weak gradient signal for the discrete node-selection head: the selection
decision influences the reward only indirectly, through the quality of the
subsequent displacement.  Possible remedies include differentiable selection
via Gumbel-Softmax, a dynamic priority queue that always selects the
highest-crossing node, or a hierarchical policy that separates region
selection from local refinement.

% ─────────────────────────────────────────────────────────────────────────────
\section{Conclusion}
\label{sec:conclusion}

We presented a systematic study of five RL formulations for graph layout
optimization, progressively addressing the limitations of each design.
The central finding is that initialization quality dominates policy
complexity: Sequential Refinement---a GAT policy that simply cycles through
nodes in BFS order starting from a \texttt{neato} layout---achieves SPC
$-24.43\%$ relative to \texttt{neato} on the Rome benchmark, far exceeding
the BFS placement approach (Method~3) and simpler neato-initialized methods.  The attempt to improve further with a learned
node-selection policy (Node-Select Refinement, SPC $-5.20\%$) highlights
the difficulty of training discrete selection decisions with sparse
RL rewards.  Future work should explore differentiable node selection and
two-stage pipelines that decouple coarse layout generation from fine-grained
crossing reduction.

% ─────────────────────────────────────────────────────────────────────────────
\bibliographystyle{plain}
\begin{thebibliography}{9}

\bibitem{schulman2017ppo}
J.~Schulman, F.~Wolski, P.~Dhariwal, A.~Radford, and O.~Klimov.
\newblock Proximal policy optimization algorithms.
\newblock \textit{arXiv preprint arXiv:1707.06347}, 2017.

\bibitem{velickovic2018gat}
P.~Veli{\v{c}}kovi{\'c}, G.~Cucurull, A.~Casanova, A.~Romero, P.~Li{\`o},
and Y.~Bengio.
\newblock Graph attention networks.
\newblock In \textit{ICLR}, 2018.

\bibitem{didimo2011rome}
W.~Didimo, G.~Liotta, and S.~A.~Romeo.
\newblock Graph drawing beyond planarity: from theory to practice.
\newblock Technical Report, 2011.

\bibitem{fruchterman1991}
T.~M.~J.~Fruchterman and E.~M.~Reingold.
\newblock Graph drawing by force-directed placement.
\newblock \textit{Software: Practice and Experience}, 21(11):1129--1164, 1991.

\bibitem{kamada1989}
T.~Kamada and S.~Kawai.
\newblock An algorithm for drawing general undirected graphs.
\newblock \textit{Information Processing Letters}, 31(1):7--15, 1989.

\end{thebibliography}

\end{document}
